\section{Introduction}

Classic commons theory predicts over-extraction under open access and highlights institutional design as the key lever for sustainable use \citep{gordon1954common,scott1955fishery,ostrom1990governing}. In machine learning, sequential social dilemma research has made the same point in a different language: cooperation is a property of policies interacting over time in stateful environments, so evaluation must account for long horizons, partial observability, and held-out social situations rather than one-shot payoffs \citep{leibo2017multiagent,meltingpot2021,du2023review,socialjax2025}. Recent LLM-agent commons work shows that cooperative behavior can emerge in simulated societies, but those findings remain fragile under changes in prompts, partner populations, and environmental conditions \citep{piatti2024cooperate,govsim_repo}. What remains weakly characterized is whether governance itself stays effective once new exploitative strategies keep entering the population.

We frame this problem as \emph{governance under adversarial strategy injection}. The central question extends beyond a single fishery result: \emph{how can governance stabilize cooperation under repeated strategic turnover in sequential commons, and how do the relevant design choices change when governance moves from fixed central signals to broader governance architectures?} To address this question, the paper is organized as two linked studies. The first study uses a minimal Fishery Commons setting to identify which \emph{central governance signals} remain effective under invasion pressure. It begins with a matched mutation-versus-live-LLM baseline, then widens into a top-down comparison over monitoring with sanctions, adaptive quotas, and temporary closures, and finally evaluates the resulting ranking with a learned-policy self-play validation in the same medium tier. The second study moves to Harvest Commons and asks a downstream question: once useful central signals have been identified, how should governance itself be organized? Earlier Harvest experiments also considered none and bottom-up-only conditions, but the decision-critical evidence in the manuscript is the narrower top-down-only versus hybrid comparison, with a learned-policy check reported only as supporting validation.

This design reflects two considerations. Commons research has long shown that communication, collective choice, and locally enforced rules can materially change outcomes even without full central command \citep{ostrom1992covenants,walker2000collective}. At the same time, benchmark work in multi-agent learning increasingly argues that robust conclusions should survive transfer across substrates and social situations, rather than depending on a single tuned environment \citep{meltingpot2021,socialjax2025}. The Fishery Commons study therefore isolates the signal-design question, whereas the Harvest Commons study widens the lens to governance architecture. Taken together, the two studies address the same overarching problem at adjacent levels: which central interventions work, and what changes when governance is no longer only a hard-coded central rule.

The paper introduces a two-level invasion-stability protocol for sequential commons. First, it uses repeated adversarial strategy turnover, held-out regime evaluation, run-level uncertainty, and learned-policy validation to compare central governance signals in a minimal renewable commons. Second, it extends the same pressure test to a new commons substrate and evaluates a narrower governance-architecture comparison between top-down-only and hybrid control, while situating that comparison within a broader experimental setting that also includes none and bottom-up-only conditions. The resulting evidence supports concrete empirical claims while remaining relevant to decentralized AI settings in which governance is part of the system design.
